\documentclass[fontsize=11pt]{article}
\usepackage{amsmath}
\usepackage[utf8]{inputenc}
\usepackage[margin=0.75in]{geometry}
\usepackage{hyperref}

\title{CSC111 Project Proposal: }
\author{Caellum Yip Hoi-Lee, Catherine Abdul-Samad, Michael Chen, Joshua Yeung}
\date{\today}

\begin{document}
\maketitle

\section*{Problem Description and Research Question}

- motivation: explain a bit of the psychology behind how obsidian tries to mimic the interwovenness of the human memory and how it's not isolated and stuff.
Personal knowledge management (PKM) tools such as Obsidian allow users to organize their knowledge as a network of interlinked notes. A note in the network can link to other notes in the network through internal links called "WikiLinks", forming a graph where the notes are vertices and the links are edges. The benefit of such a system lies in its ability to provide users insights into the relationship between ideas, and build a stronger understanding of their notes. However, many users (including ourselves), tend to under-link their notes, leading to a lackluster knowledge network. For example, one of the ways one of our group members uses Obsidian is through a "hub-and-spoke" topology, where minimal note vertices known as "hubs" form a directory with WikiLinks to other notes in the same topic (e.g. a hub named "CSC111.md" with WikiLinks to ([[Tree]], [[Recursion]], etc.). This however, defeats the purpose of Obsidian, failing to capture the lateral connections notes can have in the network which actually make a knowledge graph useful (e.g. generating links between [[Trees]] to [[Abstract Syntax Trees (CS)]] to [[Abstract Syntax Trees (Linguistics)]]). 

Whilst manually reviewing every note to add missing links between notes is a valuable learning activity due to the mental friction required, (include citation about how friction commits better to memory or something) , it becomes impractical when the PKM system grows.

\subsection{Research Question}
\textbf{"Can we use graph-based link prediction algorithms augmented with NLP to suggest meaningful missing connections in an Obsidian knowledge graph, and how do the algorithms' judgements compare with a user's own judgement of which notes should be linked?"}

\section*{Computational Plan}
\subsection*{Data Representation}
We will model an Obsidian vault as an undirected graph, $G = (V, E)$, where each Markdown note is a vertex $v \in V$ and each WikiLink between two notes is an edge $\{u, v\} \in E$. Despite graphs with WikiLinks technically being a directed graph, for example, a note titled "Graphs" would link to "Vertex" but not vice versa, we will represent the graph with an undirected graph, due to the goal being link prediction. 

Each vertex $v$ will be an object containing the following instance attributes:
- note\_name (filename minus the .md)
- note\_content (plaintext body)


We will use two datasets for our evaluation: one of our group members' personal Obsidian vaults used mainly in notetaking for school, and the open source vault of Jacky Zhao (jzhao.xyz), which contains over 740 notes spanning from computer science to philosophy. Two datasets are used due to the concern of a small dataset (mine) leading to trivial results. With a richer and denser dataset, we hope to see emergent patterns and clearer signs of effective judgement by the augmented algorithms.

Before constructing the graph, we will need to process the notes via a cleaning pipeline:
\begin{enumerate}
	\item strip image Embeddings
	\item remove tags
	\item strip frontmatter(?) / parse for metadata
	\item extract all WikiLinks
	\item extract plaintext body
\end{enumerate}

\subsection{Computations}
We plan to implement link predictions in three stages. At each stage, we will tweak a similarity signal and evaluate how the predictions change. Hopefully we will be able to deduce how different similarity signals affect the ju:dgements made by the algorithm. 

Stage 1: Structural similarity (Jaccard)
Turn Neighbourhoods into sets -> use Jaccard on sets -> get number

Stage 2: **Weighted** structural similarity (Adamic-Adar)
> A weighted version of Common Neighbours that gives higher importance to rare mutual connections
 
https://networkx.org/documentation/stable/reference/algorithms/generated/networkx.algorithms.link\_prediction.adamic\_adar\_index.html
https://www.geeksforgeeks.org/python/link-prediction-predict-edges-in-a-network-using-networkx/
https://graphty-org.github.io/algorithms/examples/algorithms/link-prediction/adamic-adar.html
cool visualisation!

Stage 3: Semantic similarity via Sentence-BERT
- sentence embeddings for each note's plaintext body with sentence-transformers lib (all-MiniLM-L6-v2?)
- use SentenceTransformer.encode() -> maps text to 384D vectors for semantic meaningful -> cosine similarity between embedding pairs
- allows us to detect similarity between notes that may share no structural similarity but have semantic similarity (related to the same thing)

Stage ? combined score:
- Compute a weighted combination of the structural score (max of Jaccard and Adamic-Adar, normalised), and semantic score

How can we quantify whether or not an evaluation is \_good\_?

Evaluation via edge holdout:
Sort of like Catherine's idea: remove a random 20\% of existing edges -> run full prediction pipeline -> measure fraction of held-out edges appearing in the top-$k$ predictions\dots

talk about the limitations of using edge holdout:
- how do you know which 20\% to remove? wouldn't randomness lead to more "important" edges being taken - 

https://www.geeksforgeeks.org/software-engineering/introduction-of-holdout-method/
https://jordan-hay.github.io/jekyll/update/data/2023/11/30/classifiereval.html

Graph statistics:
- we can also compute degree centrality per vertex to identify "hub" notes, and run connected\_components to detect isolated clusters (maybe provide info to the user on how they can improve the interwovenness of their notes)
- can compare these statistics before and after prediction pipeline to see how the predictions affect the structure of the graph

\subsection{Visualisation}
- Plotly (plotly.graph\_objects)
- NetworkX , networkx.spring\_layout


\section*{References}
\begin{enumerate}
	\item \url{https://obsidian.md/ 27 Feb. 2026.}
	\item \url{https://plotly.com/python/network-graphs/}
	\item \url{https://networkx.org/}
	\item \url{Reimers, Nils, and Iryna Gurevych. Sentence-BERT: Sentence Embeddings using Siamese BERT-Networks.}
	\item \url{Proceedings of the 2019 Conference on Empirical Methods in Natural Language Processing, 2019.}
\end{enumerate}


% NOTE: LaTeX does have a built-in way of generating references automatically,
% but it's a bit tricky to use so you are allowed to write your references manually
% using a standard academic format like MLA or IEEE.
% See project proposal handout for details.

\end{document}
